% Main document: graphicx and tikz; serif typography follows the main document.
% The same augmentations are applied to the real image and prototype thumbnails.
\begin{figure*}[t]
\centering
\begingroup
\def\MIPAImageDirectory{images}
\def\MIPACenterImage{\MIPAImageDirectory/real.png}
% Clockwise from top: filename / first-view probability / second-view probability.
% Measured softmax probabilities at tau0.025 over10 prototypes; displayed to one decimal.
\def\MIPAPrototypeData{%
  prototype-01/7.8/5.7,
  prototype-02/4.5/5.8,
  prototype-03/6.1/19.9,
  prototype-04/2.7/5.1,
  prototype-05/10.5/10.7,
  prototype-06/33.5/22.5,
  prototype-07/10.2/10.2,
  prototype-08/11.5/7.8,
  prototype-09/4.6/3.3,
  prototype-10/8.6/8.9%
}
\def\MIPAMaxPercent{35}
\definecolor{MIPAReal}{RGB}{91,132,180}
\definecolor{MIPASynthetic}{RGB}{199,106,100}
\definecolor{MIPABlueFill}{RGB}{217,230,244}
\definecolor{MIPARedFill}{RGB}{248,222,218}
\definecolor{MIPAText}{RGB}{40,44,49}
\definecolor{MIPAMuted}{RGB}{91,101,111}
\colorlet{MIPAGrid}{MIPAMuted!14!white}
\colorlet{MIPARim}{MIPAMuted!24!white}
\colorlet{MIPAFrame}{MIPAMuted!28!white}
% Both central and rim images have the same 28bp size; views are precomputed.
\newcommand{\MIPAImage}[5]{%
  \node[inner sep=0pt,outer sep=0pt] at (#2,#3)
    {\includegraphics[width=#4bp,height=#4bp]{#1}};
}

\resizebox{\linewidth}{!}{%
\begin{tikzpicture}[x=1bp,y=1bp,inner sep=0pt,outer sep=0pt,
  every node/.style={font=\rmfamily}]
  \path[use as bounding box] (0,-267) rectangle (504,0);
  \fill[white] (0,-267) rectangle (504,0);
  \foreach \MIPAView/\MIPACenterX in {1/126,2/378} {
    \edef\MIPAImageDirectory{images/view\MIPAView}
    \ifnum\MIPAView=1
      \colorlet{MIPAAccent}{MIPAReal}
      \colorlet{MIPABarFill}{MIPABlueFill}
      \def\MIPATitle{Rotation}
      \def\MIPASubtitle{$90^{\circ}$ clockwise}
    \else
      \colorlet{MIPAAccent}{MIPASynthetic}
      \colorlet{MIPABarFill}{MIPARedFill}
      \def\MIPATitle{Flip + crop}
      \def\MIPASubtitle{Horizontal flip, $1.3\!\times$ zoom}
    \fi
    \node[anchor=base,text=MIPAText,font=\rmfamily\bfseries\upshape\fontsize{10bp}{12bp}\selectfont]
      at (\MIPACenterX,-12) {\MIPATitle};
    \node[anchor=base,text=MIPAMuted,font=\rmfamily\mdseries\upshape\fontsize{8bp}{10bp}\selectfont]
      at (\MIPACenterX,-25) {\MIPASubtitle};
    \begin{scope}[shift={(\MIPACenterX,-145)}]
      % Both panels use the same radial scale.
      \foreach \MIPAFraction in {.5,1} {
        \pgfmathsetmacro{\MIPARadius}{24+39*\MIPAFraction}
        \draw[MIPAGrid,line width=.3bp,dash pattern=on .7bp off 2bp]
          (0,0) circle[radius=\MIPARadius bp];
      }
      \foreach \MIPAFile/\MIPAFirst/\MIPASecond
        [count=\MIPAIndex from 0] in \MIPAPrototypeData {
        \ifnum\MIPAView=1\let\MIPAProbability\MIPAFirst
        \else\let\MIPAProbability\MIPASecond\fi
        \pgfmathsetmacro{\MIPAAngle}{90-36*\MIPAIndex}
        \pgfmathsetmacro{\MIPAStart}{\MIPAAngle-12}
        \pgfmathsetmacro{\MIPAEnd}{\MIPAAngle+12}
        \pgfmathsetmacro{\MIPAOuter}{24+39*\MIPAProbability/\MIPAMaxPercent}
        \colorlet{MIPABar}{MIPABarFill}
        \colorlet{MIPABarEdge}{MIPAAccent!80!white}
        \draw[MIPAGrid,line width=.3bp] (\MIPAAngle:\MIPAOuter) -- (\MIPAAngle:70);
        \path[fill=MIPABar,draw=MIPABarEdge,line width=.4bp,line join=round]
          (\MIPAStart:24) -- (\MIPAStart:\MIPAOuter)
          arc[start angle=\MIPAStart,end angle=\MIPAEnd,radius=\MIPAOuter bp]
          -- (\MIPAEnd:24)
          arc[start angle=\MIPAEnd,end angle=\MIPAStart,radius=24bp] -- cycle;
      }
      \filldraw[fill=white,draw=MIPARim,line width=.4bp] (0,0) circle[radius=24bp];
      \filldraw[fill=white,draw=MIPAReal,line width=.55bp]
        (-14.7,-14.7) rectangle (14.7,14.7);
      \MIPAImage{\MIPACenterImage}{0}{0}{28}{\MIPAView}
      \foreach \MIPAFile/\MIPAFirst/\MIPASecond
        [count=\MIPAIndex from 0] in \MIPAPrototypeData {
        \ifnum\MIPAView=1\let\MIPAProbability\MIPAFirst
        \else\let\MIPAProbability\MIPASecond\fi
        \pgfmathsetmacro{\MIPAAngle}{90-36*\MIPAIndex}
        \colorlet{MIPABorder}{MIPAFrame}\colorlet{MIPALabel}{MIPAMuted}
        \pgfmathtruncatemacro{\MIPAUpper}{sin(\MIPAAngle)>0}
        \ifnum\MIPAUpper=1
          \def\MIPALabelAnchor{south}\def\MIPALabelY{17}
        \else
          \def\MIPALabelAnchor{north}\def\MIPALabelY{-17}
        \fi
        \def\MIPABorderWidth{.4bp}\let\MIPAWeight\mdseries
        \ifdim\MIPAProbability pt>20pt
          \colorlet{MIPABorder}{MIPAAccent}\colorlet{MIPALabel}{MIPAText}
          \def\MIPABorderWidth{.65bp}\let\MIPAWeight\bfseries
        \fi
        \begin{scope}[shift={(\MIPAAngle:86)}]
          \fill[white] (-16,-16) rectangle (16,16);
          \filldraw[fill=white,draw=MIPABorder,line width=\MIPABorderWidth]
            (-14.7,-14.7) rectangle (14.7,14.7);
          \MIPAImage{\MIPAImageDirectory/\MIPAFile.png}{0}{0}{28}{\MIPAView}
          \node[anchor=\MIPALabelAnchor,fill=white,inner xsep=2bp,inner ysep=1bp,
            text=MIPALabel,font=\rmfamily\MIPAWeight\upshape\fontsize{8bp}{9bp}\selectfont]
            at (0,\MIPALabelY) {\MIPAProbability\%};
        \end{scope}
      }
    \end{scope}
  }
\end{tikzpicture}%
}
\endgroup
\caption{\textbf{Soft assignments across augmentations.} A frozen TinyImageNet ResNet-18 encoder assigns the real image to the ten shown synthetic prototypes using softmax over $\ell_2$-normalized feature similarities ($\tau=0.025$). Each panel applies the same transform to all images. Probabilities are normalized over these ten prototypes; both panels use a 0--35\% radial scale.}
\label{fig:prototype-assignments}
\end{figure*}
